\documentclass{situationist-preprint}

\title{Care as the Being of Dasein: A Completeness Result}
\msid{PHE-06-a}
\msfamily{PHE-06}
\msdiscipline{Phenomenology}
\mstype{Principal manuscript}
\mskeywords{Heidegger; care; Dasein; existentials; attunement; understanding; discourse; falling; being-with; embodiment; temporality}

\begin{document}
\maketitle

\begin{abstract}
In \S41 of \emph{Being and Time} Heidegger claims that the being of Dasein is care: being ahead of itself, already in a world, as being alongside entities encountered within it. The claim has two parts. Care is supposed to unify the existentials laid out in Division~I, namely attunement, understanding, discourse and falling, and it is supposed to be complete, so that no feature of Dasein's being falls outside it. Critics from Sartre and Merleau-Ponty to Levinas have argued that the list is incomplete: being-with others, spatiality and above all embodiment are said to be further existentials that care does not contain. We defend both parts of Heidegger's claim. We state a criterion for being an existential, show that every disclosure has exactly three dimensions corresponding to the three moments of care, and show that discourse articulates the whole rather than adding a fourth moment. We then argue, case by case, that being-with, spatiality and embodiment are real and irreducible features of human existence but are ways in which the moments of care are realized, not further existentials. Finally we argue that temporality, which Heidegger calls the meaning of care, is not a further existential either: it is what makes the unity of care possible, not a member of it.
\end{abstract}

\section{The problem}

\emph{Being and Time} proceeds by an analytic of the entity that we ourselves are, which Heidegger calls Dasein. Its first division lays out a series of structures, the existentials, that characterize Dasein's way of being in contrast to the categories that characterize things \citep[\S9]{Heidegger1927}. Dasein is being-in-the-world (\S12); its world has the structure of a referential whole of significance (\S\S14--18); it is with others and is, proximally and for the most part, the anonymous \enquote{they} (\S\S25--27); it is its \enquote{there} through attunement (\emph{Befindlichkeit}), understanding (\emph{Verstehen}) and discourse (\emph{Rede}) (\S\S29--34); and in its everyday being it has fallen into the world of its concern (\S\S35--38). Chapter~6 then asks whether these structures form a whole and what its unity is. The answer is care (\emph{Sorge}): the being of Dasein is \enquote{ahead-of-itself-already-being-in-(the-world) as being-alongside (entities encountered within-the-world)} \citep[\S41]{Heidegger1927}.

The claim that care is Dasein's being contains two theses that are not always distinguished. The \emph{unity thesis} says that the existentials are not a mere aggregate but moments of one structure, none of which can occur without the others. The \emph{completeness thesis} says that the structure so unified is the whole of Dasein's being: there is no existential that is not a moment of care or a mode of one. The unity thesis has received most attention from commentators. The completeness thesis has received most attention from critics, who have argued that Division~I leaves out features of human existence that are as fundamental as any it includes. Sartre complained that Heidegger's Dasein has no body and that his being-with never reaches the concrete other \citep[Part~III, chs.~1--2]{Sartre1943}. Merleau-Ponty made the lived body, and the motor intentionality through which it inhabits space, the ground of all sense \citep{MerleauPonty1945}. Levinas argued that Dasein is never hungry, that we live from food and light and do not merely use them, and that the relation to the other person is not a structure of one's own being at all \citep{Levinas1961}. Löwith had earlier argued that Heidegger's analysis of being-with subordinates the other to the self's concern with its own being \citep{Lowith1928}, and Theunissen that it treats the other as a moment of the self's world rather than as an interlocutor \citep{Theunissen1965}. More recently \citet{Aho2009} has documented Heidegger's neglect of the body and its costs.

We argue that care is the being of Dasein in both senses: it unifies the existentials and no further existential is required. The argument has three parts. First, we state a criterion for being an existential and show that disclosedness has exactly three dimensions, which are the three moments of care; discourse articulates all three without adding a fourth (Section~\ref{s:arg}). Second, we examine the three most serious candidates for a further existential, namely being-with, spatiality and embodiment, and argue that each is a way in which the moments of care are realized, not an addition to them (Sections~\ref{s:with}--\ref{s:body}). Third, we argue that temporality, which Heidegger identifies as the meaning of care, is not a further existential but the condition of the unity of care (Section~\ref{s:temp}). The completeness thesis is compatible with conceding to the critics most of what they describe; what it denies is that what they describe must be added to care rather than found within it.

\section{Chapter six of Division I}\label{s:ch6}

The argument of \S\S39--44 is compressed and its order matters, so we set it out before stating our framework.

\paragraph{\S39: the question of totality.} Heidegger asks how the totality of the structural whole that has been exhibited is to be characterized. He rejects two answers in advance. The whole cannot be built up by piecing the elements together, since that would presuppose a plan of the whole; and it cannot be found by locating one element from which the others are derived, since the existentials are said to be equiprimordial. What is needed is a phenomenon in which the whole is disclosed at once, and in which Dasein is brought before itself as a whole.

\paragraph{\S40: anxiety.} The phenomenon is anxiety. In fear, Dasein is afraid of a definite entity within the world that threatens it. In anxiety, what one is anxious about is not any entity but being-in-the-world as such; entities within the world lose their significance, and the familiar world withdraws. Anxiety thereby individualizes Dasein, bringing it before its own possibility of being, and before the fact that it has been thrown into a world it did not choose, while showing the flight into everyday concern for what it is. Anxiety therefore discloses at once the three things Chapter~6 needs: that Dasein is concerned with its potentiality-for-being, that it already finds itself in a world, and that it is ordinarily absorbed in the world of its concern. What anxiety discloses is the topic of PHE-16, which argues that anxiety discloses the nothing as the withdrawal of beings as a whole, without revealing a further entity; we rely only on the uncontroversial point that anxiety discloses being-in-the-world as such rather than any entity within it.

\paragraph{\S41: care.} The three disclosed features are named. Dasein is \emph{ahead of itself}: in each case it is concerned with what it can be, and is its possibilities. It is \emph{already in} a world: it has been thrown and finds itself attuned to its situation. It is \emph{alongside} entities encountered within the world: it is occupied with them, and in everydayness it has fallen into this occupation. The unity of these is care. Heidegger insists that care is not a particular attitude such as worry or solicitude, and is not a disposition such as willing, wishing, urge or addiction; these are founded on care and presuppose it. Concern (\emph{Besorgen}) for things and solicitude (\emph{F\"ursorge}) for others are modes of care, not additions to it.

\paragraph{\S42: confirmation.} Heidegger offers a pre-ontological confirmation: the fable of Cura in Hyginus, in which Care shapes the human being from earth and is granted possession of it for as long as it lives. The fable is cited as testimony that Dasein's self-interpretation has long found its being in care, not as an argument.

\paragraph{\S\S43--44: reality and truth.} The final sections draw consequences. Reality is a mode of being of entities within the world and is ontologically founded in care, so that the problem of the external world, posed as a problem of proving that there are things outside a worldless subject, is a pseudo-problem (\S43). Truth in its primordial sense is the disclosedness of Dasein and the uncoveredness of entities, on which the truth of assertions is founded (\S44). These consequences show what Heidegger takes care to accomplish: it is the structure in virtue of which anything is intelligible at all.

Two features of this chapter matter for what follows. First, the three moments of care correspond to three existentials of Division~I: understanding, which projects possibilities, to being ahead of itself; attunement, which discloses thrownness, to being already in; and falling to being alongside. Discourse, which \S34 had declared equiprimordial with attunement and understanding, has no moment of its own. Second, the list is introduced as exhaustive without a separate argument. Heidegger relies on the phenomenon of anxiety to show the whole at once. Since anxiety may be held to disclose only what it discloses, the completeness of care needs an argument that does not rest on anxiety alone. We supply one.

\section{Framework}\label{s:frame}

\begin{definition}[Existential]\label{d:exist}
An \emph{existential} is a structure $S$ of Dasein's being such that (i) $S$ belongs to every possible Dasein, in every mode of its existence, and (ii) $S$ is a respect in which Dasein's being is disclosed to it, so that Dasein's understanding of its own being and of other entities depends on $S$.
\end{definition}

Clause~(i) separates existentials from features that some Dasein have and others lack: a particular mood is not an existential, but being attuned is. Clause~(ii) separates existentials from features that hold of Dasein but play no role in its disclosedness, such as having a certain mass. Clause~(ii) is Heidegger's contrast with categories (\S9): categories are determinations of entities whose being is not an issue for them, existentials are determinations of the entity whose being is an issue for it, and therefore of the ways in which it has that being disclosed \citep[\S9]{Heidegger1927}. The definition is formal: it says what would make something an existential, without assuming in advance which existentials there are.

\begin{definition}[Mode and realization]\label{d:mode}
An existential $S'$ is a \emph{mode} of an existential $S$ if every case of $S'$ is a case of $S$ and $S'$ differs from other cases of $S$ in how, not in whether, $S$ is disclosive. A structure is a \emph{realization} of $S$ for a kind of Dasein if it is the concrete way in which Dasein of that kind has $S$.
\end{definition}

Concern and solicitude are modes of being alongside in this sense; inauthenticity is a mode of each moment. The notion of realization is needed for the critics' cases. If every human Dasein is attuned through bodily feeling, bodily feeling is the realization of attunement for human Dasein. A realization is not a further existential, because it does not add a respect in which Dasein's being is disclosed; it is the way in which an existing respect is had.

\begin{definition}[Unity and completeness]\label{d:comp}
A set $C$ of existentials is \emph{unified} if no member of $C$ can be had by a Dasein without the others. $C$ is \emph{complete} if every existential is a member of $C$ or a mode of a member of $C$. Care is the being of Dasein if the moments of care are unified and complete.
\end{definition}

The completeness thesis is thus a universal claim about existentials, and a candidate counterexample must satisfy both clauses of Definition~\ref{d:exist}. A critic who shows that embodiment belongs to every human being but is not disclosive in a distinct respect has not produced a counterexample; nor has one who shows that it is disclosive but only for some Dasein.

\begin{principle}[Disclosure is mattering]\label{p:matter}
Every disclosure is a disclosure of something as mattering in some way, and something can matter only (a) in the light of some possibility of the one to whom it matters, (b) in some respect in which that one already finds itself affected, and (c) as something with which that one is occupied or is in a position to be occupied.
\end{principle}

Principle~\ref{p:matter} is the substantive premise of the argument. Its first part is Heidegger's thesis that Dasein encounters entities in their significance, not first as bare things to which value is later added (\S\S15--18), a thesis defended at length by \citet{Dreyfus1991} and \citet{Haugeland1982}. Its second part says what mattering requires. Each clause can be motivated independently. Without (a), nothing would count as relevant to anything, since relevance is relevance to what one can do or be: a hammer is disclosed as too heavy only in the light of hammering, and a remark as insulting only in the light of the one it insults being someone who can be insulted. Without (b), nothing could affect the one to whom it is disclosed, since a thing that matters must already be able to touch one; it is through being attuned that a situation is disclosed as threatening, promising or dull, as \S29 argues and as PHE-21 develops in arguing that moods are pre-reflective disclosures of one's situation, fallible but not arbitrary. Without (c), there would be no \emph{something} that matters, but at most a possibility and an affection without an object. The three clauses are thus not three things that happen to accompany every disclosure but three places in the structure of mattering, analogous to the three places of givenness defended in PHE-02, on which something is given to someone as something.

\begin{principle}[Articulation is not a further place]\label{p:artic}
That a disclosure is articulated, so that what is disclosed has parts and joints that can be expressed, is a feature of the whole disclosure and not a further place in its structure.
\end{principle}

Principle~\ref{p:artic} concerns discourse. In \S34 Heidegger defines discourse as the articulation of intelligibility: it is what makes what is understood and attuned capable of being laid out in its joints, prior to its being put into words \citep[\S34]{Heidegger1927}. \citet{Carman2003} and \citet{Wrathall2010} have argued that discourse in this sense is neither language nor communication but the structured intelligibility of what is disclosed, which language expresses. The principle says that such articulation is a character of the whole, as the grammatical structure of a sentence is not a further word in it.

\section{The argument}\label{s:arg}

\begin{lemma}[Three dimensions]\label{l:three}
Disclosedness has exactly three dimensions: being ahead of oneself in possibilities, being already in a situation that affects one, and being alongside what is encountered.
\end{lemma}

\begin{proof}
By Definition~\ref{d:exist}, existentials are respects in which Dasein's being is disclosed. By Principle~\ref{p:matter}, every disclosure involves something mattering in the light of a possibility, in an affected respect, and as something encountered. These are three places in a single relation, and each is required. That there are no more places follows from the principle's \enquote{only}: anything further that a disclosure involves is either a further determination of one of the three (which possibility, which affection, which entity) or a feature of the whole, such as its articulation, which by Principle~\ref{p:artic} is not a further place. The three places are what Heidegger calls being ahead of oneself, being already in, and being alongside \citep[\S41]{Heidegger1927}.
\end{proof}

\begin{lemma}[Correspondence]\label{l:corr}
Understanding is the existential of the first dimension, attunement of the second and falling, as the everyday mode of being alongside, of the third. Discourse is the articulation of all three.
\end{lemma}

\begin{proof}
Understanding in Heidegger's sense is not cognition but Dasein's being its possibilities, the projection of its being upon the for-the-sake-of-which and upon significance (\S31). It is therefore the existential through which Dasein is ahead of itself. Attunement discloses Dasein's thrownness, the fact that it is and has to be, and makes it possible to be affected by entities (\S29); it is therefore the existential through which Dasein is already in a world. Being alongside entities is, in everydayness, falling: absorption in the world of concern and in the public interpretedness of the \enquote{they}, articulated as idle talk, curiosity and ambiguity (\S\S35--38). Falling is a mode of being alongside in the sense of Definition~\ref{d:mode}, since it is the way in which being alongside is had proximally and for the most part, not a further respect of disclosure. Discourse is equiprimordial with attunement and understanding (\S34) and is the articulation of intelligibility; by Principle~\ref{p:artic} it is a character of the whole disclosure, and so it articulates each dimension rather than constituting a further one.
\end{proof}

Lemma~\ref{l:corr} explains a feature of the text that has puzzled commentators. Discourse has no place in the formula of care, and in \S68(d) Heidegger says that discourse does not temporalize itself primarily in any one ecstasis of temporality \citep[\S68]{Heidegger1927}. If discourse were a moment like the others, this would be an anomaly. On our reading it is what one should expect: discourse articulates the whole, so it belongs to no single moment and has no single ecstasis of its own. \citet{Dreyfus1991} treats the place of discourse as a sign of unclarity in Heidegger's account; we take it to be a consequence of what discourse is.

\begin{lemma}[Unity]\label{l:unity}
The three moments of care are unified: none can be had without the others.
\end{lemma}

\begin{proof}
Suppose Dasein were ahead of itself without being already in a world. Its possibilities would then be possibilities of nothing in particular, since a possibility of Dasein is always a possibility of acting in a situation, with these resources and against these constraints; possibilities are disclosed only as possibilities of the thrown. Suppose it were already in a world without being ahead of itself. Its situation would not be disclosed as a situation, since a situation is disclosed only in the light of what one can do in it; attunement without projection would affect nothing. Suppose it were ahead of itself and already in a world but alongside nothing. Its possibilities would be possibilities of dealing with nothing, and its attunement an affection by nothing, contrary to Principle~\ref{p:matter}(c). Conversely, being alongside entities without projection or attunement would be the mere co-presence of things, which is not a way of being disclosed at all. Each moment therefore requires the others.
\end{proof}

\begin{theorem}[Completeness]\label{t:comp}
Every existential is a moment of care or a mode of one. Hence care is the being of Dasein.
\end{theorem}

\begin{proof}
Let $S$ be an existential. By Definition~\ref{d:exist}(ii), $S$ is a respect in which Dasein's being is disclosed. By Lemma~\ref{l:three}, every respect of disclosure falls within one of the three dimensions or is the articulation of the whole. In the first case $S$ is either the existential of that dimension or a mode of it, by Lemma~\ref{l:corr}; in the second case $S$ is discourse, which Lemma~\ref{l:corr} shows articulates the moments of care and is equiprimordial with them. So the moments of care, with their articulation, are complete in the sense of Definition~\ref{d:comp}. By Lemma~\ref{l:unity} they are unified. By Definition~\ref{d:comp}, care is the being of Dasein.
\end{proof}

The proof shows where the critics must attack. They must deny Principle~\ref{p:matter} or Principle~\ref{p:artic}, or they must produce a structure that satisfies Definition~\ref{d:exist} and is not a mode of a moment of care. The second line is the one they have taken. The next three sections consider the candidates.

\section{Being-with}\label{s:with}

Heidegger himself calls being-with (\emph{Mitsein}) an existential and says that Dasein is essentially being-with even when no other is present or perceived, since being alone is a deficient mode of being-with (\S26). Being-with is not mentioned in the formula of care. One might therefore conclude that Heidegger's own list exceeds care, and that the completeness thesis is false by his own lights. This is the internal form of the objection. Its external form, due to \citet{Sartre1943} and \citet{Levinas1961}, says that Heidegger's being-with is the wrong structure for the relation to the other and that the right one is not a structure of care at all.

The internal objection mistakes the place of being-with. It is not a fourth moment alongside the other three, since it is not a further respect of disclosure: it is a character of each. Dasein's possibilities are disclosed as possibilities that others have or could have, and its for-the-sake-of-which is drawn from the public interpretedness of what one does (\S27). Its thrownness is thrownness into a world already shared, with a language, practices and roles that others have shaped. What it is alongside includes others, encountered not as present-at-hand things nor as equipment but as Dasein-with, with whom one is in solicitude (\S26). Being-with is thus a feature of the world in which Dasein is, and the world is a constituent of the care structure, since care is being-in-the-world \citep[\S41]{Heidegger1927}. In the terms of Definition~\ref{d:mode}, solicitude is the mode of being alongside whose object is other Dasein; being-with is the sharedness that characterizes all three moments. This is why Heidegger can say both that being-with is an existential and that care is the whole: being-with is an existential that belongs to care's structure, as the world does, not one that stands beside it.

The external objection grants this and complains of it. Sartre argues that Heidegger's being-with is a crew-like \enquote{we} that tells us nothing about the concrete encounter with this other, whose look makes me an object \citep[Part~III, ch.~1]{Sartre1943}; PHE-09 examines that analysis and argues that it establishes objectification, not hostility. Levinas argues that the other person is not disclosed within any horizon of mine but faces me with a demand that ruptures my world \citep{Levinas1961}. Theunissen argues that a transcendental account of the other, Husserlian or Heideggerian, cannot reach the dialogical relation to a second person \citep{Theunissen1965}.

We grant that the concrete encounter has features that Heidegger's account of solicitude does not describe. But these features do not satisfy Definition~\ref{d:exist}. Consider the Levinasian case, which is the strongest. The demand of the face is, on Levinas's account, a demand that I receive and to which I am answerable; PHE-13 argues that it is infinite in principle and cannot be fully discharged without losing its normative force. Being answerable to a demand is a way in which my possibilities are disclosed to me: it is being ahead of myself in a possibility, that of responding, that I did not project and cannot fully realize. It is also a way of being affected, an attunement, and it concerns someone I am alongside. The Levinasian demand therefore falls within the three dimensions, as a mode of the first in which the possibility is given to me rather than projected by me. \citet{Crowell2013} has argued that care in Heidegger is already a structure of answerability, in which Dasein is responsive to the normative claims of what it encounters and is called to be a ground of its own reasons; on that reading the face is a privileged instance of what care makes possible, not an exception to it. Whether the face has the ethical priority Levinas claims is a separate question, which the completeness thesis does not settle. The thesis concerns the structure of Dasein's being, not the order of ethics and ontology.

\section{Spatiality}\label{s:space}

Dasein is spatial: it is here, things are near or far from it, and its world has regions in which things belong (\S\S22--24). Heidegger describes Dasein's spatiality by two characters, de-severance (\emph{Ent-fernung}), the bringing-near of what is far, and directionality (\emph{Ausrichtung}), the orientation by which things are to the right or left, before or behind. He says that Dasein's spatiality is grounded in being-in-the-world and does not ground it (\S24), and in \S70 he attempts to derive it from temporality. In \enquote{Time and Being} he later declared that derivation untenable \citep{Heidegger1969}. The retraction has encouraged the view that spatiality is an independent existential that care, unfolded temporally, cannot contain \citep{Malpas2006}.

We distinguish the claim that spatiality cannot be derived from temporality from the claim that it is not a mode of care. The first may be true; the second does not follow. De-severance is a way of being alongside entities: to bring something near is to be occupied with it, and the nearness that de-severance establishes is measured by concern, not by distance, so that the friend I am waiting for at the end of the street is nearer than the spectacles on my nose (\S23). Directionality is a way of being already in a world: Dasein finds itself oriented, with a right and a left that it did not choose and that structure how things can affect it. Regions are articulations of the significance of the world, and are disclosed in the light of possibilities: the kitchen is where the things for cooking belong. Spatiality thus satisfies clause~(i) of Definition~\ref{d:exist}, since every Dasein is spatial, but it does not add a respect of disclosure beyond the three; it is the way in which all three are had by a being that is here rather than there.

\citet{Dreyfus1991} argued that Heidegger's account confuses Dasein's own spatiality, the nearness of concern, with the public space in which equipment is located, and that the second cannot be derived from the first. \citet{Arisaka1995} replied that the two are distinguished in the text and that public space is founded on the shared character of de-severance. Either way, public space is a feature of the world, and the world, as Section~\ref{s:with} noted, is a constituent of care. The debate concerns how spatiality is founded, not whether it is a fourth moment. Malpas's stronger thesis, that place is the condition of the disclosure of being, belongs to an interpretation of the later Heidegger, and if correct would show that the meaning of care is spatio-temporal rather than temporal. That is a claim about Section~\ref{s:temp}'s question, not about completeness.

\section{Embodiment}\label{s:body}

The most serious candidate is the body. Heidegger says in \S23 that Dasein's bodily nature hides a whole problematic of its own, which is not to be treated there \citep[\S23]{Heidegger1927}. In the Zollikon seminars, asked about Sartre's complaint that \emph{Being and Time} gives the body only a few lines, he replied that the bodily is the most difficult problem and that he had not been able to say more about it at the time \citep{Heidegger1987}. Three critics have argued that the omission is not a gap in exposition but a gap in the ontology.

\paragraph{Sartre.} In \emph{Being and Nothingness}, Sartre distinguishes three ontological dimensions of the body: the body as I exist it, which is not an object for me but the point of view I am; the body as known and used by the other; and the body as I exist it as known by the other \citep[Part~III, ch.~2]{Sartre1943}. He argues that the first dimension is the contingent necessity of my being in a situation: I must have a point of view, and my body is that I have this one.

\paragraph{Merleau-Ponty.} In the \emph{Phenomenology of Perception}, the lived body is the subject of a motor intentionality that is neither thought nor reflex: I reach for the glass without representing its position, and the habitual body inhabits space and tools so that they become extensions of it \citep{MerleauPonty1945}. On Merleau-Ponty's view the body is our anchorage in the world, the vehicle of being-in-the-world, and perception, which it makes possible, is the founding layer of all sense.

\paragraph{Levinas.} In \emph{Totality and Infinity}, Levinas argues that we live from good soup, air, light and sleep, which are not equipment for the sake of anything but objects of enjoyment, and that Heidegger's analysis, which encounters things as ready-to-hand within a totality of involvements, cannot describe this. Dasein, he remarks, is never hungry \citep[Section~II]{Levinas1961}.

\citet{Aho2009} argues that these complaints jointly show that the bodily is a dimension of existence that is neither reducible to the existentials of Division~I nor merely ontic, since the lived body belongs to the way the world is disclosed. We grant the second point and deny the first.

That the lived body belongs to disclosure is correct. PHE-08 argues that the lived body is the zero-point of orientation and cannot be fully given to itself as an object; it is therefore not a thing within the world but a condition of the world's being disclosed from somewhere. But to be a condition of disclosure in this sense is to be the realization, for human Dasein, of being already in a world. Sartre's analysis says as much: the body as lived is the contingent necessity of having a point of view, that is, thrownness as it is had by a being with sense organs. Merleau-Ponty's motor intentionality is the way in which human Dasein is ahead of itself in possibilities of action: the habitual body projects a practical field of what can be reached, grasped and walked across, which is understanding in Heidegger's sense, realized in a skilled organism. Levinas's enjoyment is a way of being alongside entities that is not equipmental; it shows that Heidegger's account of being alongside is too narrow when it is confined to concern with the ready-to-hand, and that it needs a mode in which things are lived from rather than used. It does not show that enjoyment is a fourth respect of disclosure, since what is enjoyed matters in the light of a possibility (of being nourished, warmed, rested), in a respect in which one is affected (hunger, cold, fatigue) and as something one is occupied with. The three places of Principle~\ref{p:matter} are all filled.

The embodied features thus fail clause~(ii) of Definition~\ref{d:exist} as further existentials: they are disclosive, but not in a further respect. They are realizations of the moments of care in the sense of Definition~\ref{d:mode}. This verdict is not a demotion. Realizations are not optional or superficial; the realization of care for human Dasein is all of human existence. What the verdict denies is only that the care structure must be supplemented. It also yields a reading of Heidegger's reticence: if the body is the realization of care for human Dasein, an ontology of Dasein that abstracts from what kind of entity realizes it will say little about the body, and the omission is a limitation of the level of description, not of the inventory of existentials. One may regret the abstraction and still accept the inventory.

\citet{Carman2003} and \citet{Wrathall2010} have both stressed that Heidegger's understanding is primarily practical, a know-how rather than a knowing-that, and \citet{Dreyfus1991} made skilled coping the model for all of Division~I. On these readings Merleau-Ponty's motor intentionality is not a rival to understanding but the most developed account of it.

\section{Temporality as the meaning of care}\label{s:temp}

Division~II argues that the unity of care is made possible by temporality: being ahead of itself is grounded in the future, being already in in having-been, and being alongside in the present, and these three ecstases form a unity in which the future has priority \citep[\S65]{Heidegger1927}. Heidegger calls temporality the \emph{meaning} of care. It may seem that this undoes the completeness thesis: if temporality grounds care, then care is not the whole of Dasein's being, and temporality is a further existential beneath it. We argue that it is not.

\begin{definition}[Meaning]\label{d:meaning}
The \emph{meaning} of a structure $S$ is that upon which $S$ is projected in being understood, so that $S$ is intelligible as the unity it is \citep[\S\S32, 65]{Heidegger1927}.
\end{definition}

\begin{proposition}\label{p:temp}
Temporality is not an existential additional to the moments of care.
\end{proposition}

\begin{proof}
Suppose temporality were a further existential $T$ of the same order as the moments of care. Then the set of existentials would be care's three moments plus $T$, and by Definition~\ref{d:comp} the question of their unity would arise again: something would have to make the four a whole. Either that unity is a fifth existential, and the regress continues, or it is not an existential but the unity of the whole. The second option is the one Heidegger takes, and it is the one Definition~\ref{d:meaning} describes: temporality is not a member of the structure but that which makes it intelligible as a unity. Moreover, temporality does not satisfy clause~(ii) of Definition~\ref{d:exist} as a separate respect of disclosure. It is not something disclosed alongside possibilities, situations and entities but the way each of these is disclosed: possibilities as coming towards Dasein, its situation as what it has been, entities as what it makes present. Since each ecstasis is the temporal character of one moment, and the moments exhaust disclosure (Lemma~\ref{l:three}), temporality adds no further respect. Hence $T$ is not an existential additional to the moments of care.
\end{proof}

The proposition reconciles the two characterizations Heidegger gives. Care is the being of Dasein, in the sense that its moments are the complete and unified articulation of how Dasein is. Temporality is the meaning of that being, in the sense that it is why these moments hang together. The relation is not one of a further layer beneath care but of the unity of a structure to the structure. This is why Heidegger can say that temporality \enquote{temporalizes} rather than is, and why it is not one more feature listed among others \citep[\S65]{Heidegger1927}.

Proposition~\ref{p:temp} does not commit us to Heidegger's further claims about temporality. \citet{Blattner1999} has argued that Heidegger's attempt to derive ordinary time from originary temporality fails, and that his position is a form of temporal idealism on which time depends on Dasein. Neither the success of that derivation nor the idealism is needed here. We need only that temporality is the condition of the unity of care rather than a member of it, which Blattner's reconstruction of originary temporality as the structure of Dasein's projection supports.

The proposition also locates the problem of wholeness with which Division~II begins. Since Dasein is always ahead of itself, it seems never to be whole while it exists; the problem is solved by being-towards-death, the possibility of no longer being able to be there, which Dasein can anticipate. PHE-07 argues that authentic being-towards-death is possible as an orientation but not as a stable achievement, since anticipation is constantly drawn back into the public time of the \enquote{they}. The completeness thesis is compatible with this. Being-towards-death is the extreme mode of being ahead of oneself, not a fourth moment, and if authenticity is only an orientation, it is an orientation of care, not a state that completes it.

\section{Recent interpretations}\label{s:lit}

The completeness thesis can be tested against the main recent readings of Division~I, each of which offers its own account of what care unifies.

\citet{Dreyfus1991} reads Division~I as an account of absorbed coping, the skilled, non-representational engagement with equipment and with others that underlies all explicit cognition. Care is the structure of such coping: projecting into possibilities, already attuned to what matters, absorbed in the task. Dreyfus's reading makes the unity thesis easy, since coping visibly involves all three moments at once. It makes the completeness thesis harder in one respect, since Dreyfus treats falling as the conformism of the \enquote{they}, and conformism seems a contingent feature of social life rather than an existential. Our reading separates the two: falling, as the everyday mode of being alongside, is the absorption in what one is occupied with; that this absorption is shaped by the public interpretedness of the \enquote{they} is the shared character of the world (Section~\ref{s:with}), not a further existential.

\citet{Blattner2006} reads care as a structure of self-interpretation: Dasein is the entity whose being is at issue for it, and care articulates how it takes up its being. On this reading the moments of care are, roughly, what one is projecting (one's self-understanding), what one has to work with (thrownness), and what one deals with in pursuing it. Blattner's emphasis on self-interpretation fits Principle~\ref{p:matter}(a) well and helps with the case of being-with, since self-interpretation is always in terms drawn from a shared world.

\citet{Haugeland1982} argued that Dasein is not a human being but a way of life, a pattern of normatively governed conduct in which individual people participate; in later work he read care as existential commitment, the resolute holding of oneself to the standards that make entities intelligible \citep{Haugeland2013}. Haugeland's reading raises the question whether care is the being of an individual or of a community. The completeness result is neutral: it concerns the structure of whatever entity has its being as an issue, and Haugeland's Dasein, if it is such an entity, has the three moments.

\citet{Crowell2013} reads care as the structure of a being that is answerable to norms, and conscience, in Division~II, as the call to be a ground of the reasons on which one acts. His reading supports our reply to Levinas: a being whose being is care is already exposed to claims, and the face is one such claim. Crowell also shows why the completeness thesis matters for Heidegger's project: if care were incomplete, there would be a respect of Dasein's being that was not answerable, and the account of conscience would lack its full scope.

\citet{Carman2003} stresses the role of discourse and interpretation, arguing that Heidegger's hermeneutic account of intelligibility is best understood in terms of the expressive articulation of what is disclosed. \citet{Wrathall2010} reads discourse as the articulation of the world into meaningful joints, independent of and prior to language. Both readings support Principle~\ref{p:artic}: discourse is how what is disclosed is structured, not one more thing disclosed. They also explain why discourse, though equiprimordial, has no moment in the formula of care.

Finally, \citet{Ratcliffe2008} and \citet{ElpidorouFreeman2015} have developed the account of attunement into an account of existential feelings, background orientations that structure the space of possibilities in which particular emotions occur. These accounts strengthen Lemma~\ref{l:unity}: an existential feeling is at once an attunement and a shaping of possibilities, which is what the unity of the first two moments predicts.

\section{Objections and replies}\label{s:obj}

\begin{objection}[Circularity]\label{o:circ}
Principle~\ref{p:matter} is the care structure in another vocabulary. The completeness result is therefore assumed rather than proved.
\end{objection}

\begin{reply}
The principle is stated in terms of mattering, which is a phenomenon available to description independently of Heidegger's terminology, and each of its three clauses was motivated by showing what would be lost without it. That the three clauses turn out to correspond to the three moments of care is the substance of the result, as the correspondence between the three places of givenness and the structure of intentionality is the substance of PHE-02. A circular argument would assume that every existential is a moment of care; ours assumes only that every existential is disclosive and that disclosure has a certain structure. A critic who denies the latter owes a case of disclosure in which something matters without mattering in the light of a possibility, in an affected respect, or as something encountered. The candidates of Sections~\ref{s:with}--\ref{s:body} were the best such cases, and each filled all three places.
\end{reply}

\begin{objection}[Realization is a further existential]\label{o:real}
If every human Dasein is embodied and embodiment is disclosive, then embodiment satisfies both clauses of Definition~\ref{d:exist}. Calling it a realization is a verbal manoeuvre.
\end{objection}

\begin{reply}
Clause~(i) of Definition~\ref{d:exist} requires that the structure belong to every possible Dasein, not every human one. The question whether there could be a Dasein that is not embodied in the human way, or not embodied at all, is open, and Heidegger's analysis abstracts from it; PID-10 argues that disembodied persons are possible in principle, and if some such person were a Dasein, embodiment would fail clause~(i). But suppose that every possible Dasein is embodied. Then embodiment satisfies clause~(i). It still does not satisfy clause~(ii) as an additional existential, because the respects in which it is disclosive are the three of Lemma~\ref{l:three}, as Section~\ref{s:body} showed case by case. The distinction between a further respect of disclosure and a further way of having the same respects is not verbal; it is the distinction between adding a dimension to a space and specifying how the space is occupied.
\end{reply}

\begin{objection}[Heidegger's own list]\label{o:list}
Heidegger lists being-with and being-one's-self as existentials in \S\S25--27, and spatiality in \S\S22--24. Completeness contradicts his own text.
\end{objection}

\begin{reply}
Heidegger's claim in \S41 is that care is the being of Dasein, which includes the existentials discussed earlier. The existentials of \S\S22--27 characterize the world and the who of being-in-the-world; the moments of care articulate being-in-the-world as a whole. There is no contradiction in holding that being-with and spatiality are existentials and that they belong to the care structure as characters of its moments. Definition~\ref{d:comp} makes this precise: completeness requires that every existential be a moment of care or a mode of one, not that only three existentials exist.
\end{reply}

\begin{objection}[Indifference]\label{o:indiff}
There are states in which nothing matters: profound indifference, apathy, depersonalization. If care is the being of Dasein, such states would be impossible, but they occur.
\end{objection}

\begin{reply}
That nothing matters to someone is itself a way in which her situation is disclosed: as flat, as lacking possibilities, as not touching her. It is an attunement, a privative mode of being ahead of oneself in which possibilities show up as unavailable, and a mode of being alongside in which entities are encountered as indifferent. Anxiety, on which Heidegger's chapter turns, is a case in which the significance of the world collapses, and it is the phenomenon that discloses care as a whole. Privative modes confirm the structure they deprive. PHE-04 makes the parallel point that profound boredom, the apparent exception to intentionality, is intentional towards time itself.
\end{reply}

\begin{objection}[The other is not mine]\label{o:levinas}
Locating the face within care, as a mode of being ahead of oneself, is exactly the reduction of the other to the same that Levinas denounced.
\end{objection}

\begin{reply}
The reply of Section~\ref{s:with} does not say that the other is constituted by my care or that the demand originates in my projection. It says that my being answerable to a demand that does not originate in me is a way in which my possibilities are disclosed to me. That the demand comes from outside, and exceeds anything I can project, is a feature of what is disclosed, which our account preserves. What Levinas denies, that the relation to the other is a structure of my being, is a claim about the source and priority of the demand; the completeness thesis concerns only how its reception belongs to the being that receives it.
\end{reply}

\section{What the result does not establish}\label{s:not}

The result is a claim about the structure of Dasein's being, and it leaves several questions open. It does not establish that Heidegger's ontology is correct against naturalistic accounts of the mind; it shows that if Dasein's being is to be described in existentials, care is the complete description. It does not establish Heidegger's derivation of spatiality or of ordinary time from originary temporality, both of which we have set aside. It does not settle the priority of ethics over ontology, which is a question about the demand of the other rather than about the being that receives it. It does not establish that the realizations of care are uninteresting: an ontology of human existence that stopped at care would be incomplete as an account of human beings, though complete as an account of Dasein's being. And it does not show that authenticity is attainable as a stable state, a question PHE-07 answers in the negative.

\section{Conclusion}

Care is the being of Dasein. Its three moments, being ahead of oneself, being already in a world and being alongside entities, correspond to the three places in the structure of anything's mattering, and discourse articulates all three without adding a fourth. The critics' candidates for further existentials are real features of human existence: being-with is the sharedness of each moment, spatiality the way all three are had by a being that is here, and embodiment the realization of the whole structure in a living organism. None adds a respect of disclosure. Temporality is not a further existential either, but the meaning of care, the unity that makes its moments a whole. The completeness thesis can therefore grant the phenomenological descriptions of Sartre, Merleau-Ponty and Levinas while denying the ontological conclusion they drew from them.

\begin{thebibliography}{99}\small
\bibitem[Aho(2009)]{Aho2009} Aho, K.~A. (2009). \emph{Heidegger's Neglect of the Body}. State University of New York Press.
\bibitem[Arisaka(1995)]{Arisaka1995} Arisaka, Y. (1995). On Heidegger's theory of space: A critique of Dreyfus. \emph{Inquiry}, 38(4), 455--467.
\bibitem[Blattner(1999)]{Blattner1999} Blattner, W.~D. (1999). \emph{Heidegger's Temporal Idealism}. Cambridge University Press.
\bibitem[Blattner(2006)]{Blattner2006} Blattner, W. (2006). \emph{Heidegger's Being and Time: A Reader's Guide}. Continuum.
\bibitem[Carman(2003)]{Carman2003} Carman, T. (2003). \emph{Heidegger's Analytic: Interpretation, Discourse, and Authenticity in Being and Time}. Cambridge University Press.
\bibitem[Crowell(2013)]{Crowell2013} Crowell, S. (2013). \emph{Normativity and Phenomenology in Husserl and Heidegger}. Cambridge University Press.
\bibitem[Dreyfus(1991)]{Dreyfus1991} Dreyfus, H.~L. (1991). \emph{Being-in-the-World: A Commentary on Heidegger's Being and Time, Division I}. MIT Press.
\bibitem[Elpidorou and Freeman(2015)]{ElpidorouFreeman2015} Elpidorou, A. and Freeman, L. (2015). Affectivity in Heidegger I: Moods and emotions in \emph{Being and Time}. \emph{Philosophy Compass}, 10(10), 661--671.
\bibitem[Haugeland(1982)]{Haugeland1982} Haugeland, J. (1982). Heidegger on being a person. \emph{No\^us}, 16(1), 15--26.
\bibitem[Haugeland(2013)]{Haugeland2013} Haugeland, J. (2013). \emph{Dasein Disclosed: John Haugeland's Heidegger}. Ed. J. Rouse. Harvard University Press.
\bibitem[Heidegger(1927)]{Heidegger1927} Heidegger, M. (1927). \emph{Sein und Zeit}. Max Niemeyer. Trans. J. Macquarrie and E. Robinson, \emph{Being and Time}, Harper \& Row, 1962.
\bibitem[Heidegger(1969)]{Heidegger1969} Heidegger, M. (1969). \emph{Zur Sache des Denkens}. Max Niemeyer. Trans. J. Stambaugh, \emph{On Time and Being}, Harper \& Row, 1972.
\bibitem[Heidegger(1987)]{Heidegger1987} Heidegger, M. (1987). \emph{Zollikoner Seminare}. Ed. M. Boss. Vittorio Klostermann. Trans. F. Mayr and R. Askay, \emph{Zollikon Seminars}, Northwestern University Press, 2001.
\bibitem[Levinas(1961)]{Levinas1961} Levinas, E. (1961). \emph{Totalit\'e et infini: Essai sur l'ext\'eriorit\'e}. Martinus Nijhoff. Trans. A. Lingis, \emph{Totality and Infinity}, Duquesne University Press, 1969.
\bibitem[L\"owith(1928)]{Lowith1928} L\"owith, K. (1928). \emph{Das Individuum in der Rolle des Mitmenschen}. Drei Masken Verlag.
\bibitem[Malpas(2006)]{Malpas2006} Malpas, J. (2006). \emph{Heidegger's Topology: Being, Place, World}. MIT Press.
\bibitem[Merleau-Ponty(1945)]{MerleauPonty1945} Merleau-Ponty, M. (1945). \emph{Ph\'enom\'enologie de la perception}. Gallimard. Trans. D.~A. Landes, \emph{Phenomenology of Perception}, Routledge, 2012.
\bibitem[Ratcliffe(2008)]{Ratcliffe2008} Ratcliffe, M. (2008). \emph{Feelings of Being: Phenomenology, Psychiatry and the Sense of Reality}. Oxford University Press.
\bibitem[Sartre(1943)]{Sartre1943} Sartre, J.-P. (1943). \emph{L'\^etre et le n\'eant}. Gallimard. Trans. H.~E. Barnes, \emph{Being and Nothingness}, Philosophical Library, 1956.
\bibitem[Theunissen(1965)]{Theunissen1965} Theunissen, M. (1965). \emph{Der Andere: Studien zur Sozialontologie der Gegenwart}. de Gruyter.
\bibitem[Wrathall(2010)]{Wrathall2010} Wrathall, M.~A. (2010). \emph{Heidegger and Unconcealment: Truth, Language, and History}. Cambridge University Press.
\end{thebibliography}

\end{document}
